\chapter{LCS con un numero di simboli rossi multiplo di 3}

\section{Definizione del problema}

\subsection{Istanza}
Siano date due sequenze di elementi sull'alfabeto $\Sigma$:
\begin{align*}
    X &= \langle x_1, x_2, \dots, x_m \rangle \\
    Y &= \langle y_1, y_2, \dots, y_n \rangle
\end{align*}
e una funzione di colorazione $\text{col} : \Sigma \to \{\text{rosso}, \text{blu}, \text{nero}\}$.

\subsection{Soluzione}
Una più lunga sottosequenza comune tra $X$ e $Y$ la quale contiene un numero di elementi rossi che è multiplo di $3$.

\section{Definizione del sottoproblema e del coefficiente}
Per ogni $i \in \{0, \dots, m\}$, per ogni $j \in \{0, \dots, n\}$ e per ogni $p \in \{0, 1, 2\}$, siano $X_i$ e $Y_j$ i prefissi di lunghezza $i$ e $j$ di $X$ e $Y$, e sia $p$ il resto della divisione per $3$ del numero $k$ di simboli rossi contenuti nella sottosequenza ($k \bmod 3 = p$).

\medskip
\noindent
\textbf{Definizione del coefficiente $c_{i,j}^p$:} Per ogni $i \in \{0, \dots, m\}$, $j \in \{0, \dots, n\}$ e $p \in \{0, 1, 2\}$, $c_{i,j}^p$ è la lunghezza della più lunga sottosequenza comune tra $X_i$ e $Y_j$ il cui numero di elementi rossi $k$ soddisfa $k \equiv p \pmod 3$:
\begin{itemize}
    \item $p = 0$: il numero di elementi rossi è un multiplo di $3$ ($k \bmod 3 = 0$);
    \item $p = 1$: il numero di elementi rossi diviso per $3$ dà resto $1$ ($k \bmod 3 = 1$);
    \item $p = 2$: il numero di elementi rossi diviso per $3$ dà resto $2$ ($k \bmod 3 = 2$).
\end{itemize}
Se per una tripla $(i, j, p)$ non esiste alcuna sottosequenza comune ammissibile, si pone $c_{i,j}^p = -\infty$.

\section{Equazioni di ricorrenza}

\subsection{Caso base}
\begin{align*}
    (i = 0 \lor j = 0) \land p = 0 &\implies c_{i,j}^0 = 0 \\
    (i = 0 \lor j = 0) \land p \in \{1, 2\} &\implies c_{i,j}^p = -\infty
\end{align*}
\textbf{Giustificazione:} La sottosequenza comune a un prefisso vuoto è $\langle \rangle$, che contiene $0$ simboli rossi. Poiché $0 \bmod 3 = 0$, per $p = 0$ si ha lunghezza $0$, mentre per $p \in \{1, 2\}$ è impossibile formare sottosequenze dal vuoto ($-\infty$).

\subsection{Passo ricorsivo}
Per $i > 0 \land j > 0$ e per ogni $p \in \{0, 1, 2\}$:
\begin{itemize}
    \item \textbf{Caso $x_i \neq y_j$:}
    \[
        c_{i,j}^p = \max \{ c_{i-1, j}^p, \, c_{i, j-1}^p \}
    \]
    \textbf{Giustificazione:} Gli elementi correnti non coincidono; il resto modulo $3$ rimane invariato e si eredita il massimo tra i sottoproblemi $(i-1, j, p)$ e $(i, j-1, p)$.

    \item \textbf{Caso $x_i = y_j$ con $\text{col}(x_i) \neq \text{rosso}$:}
    \[
        c_{i,j}^p = c_{i-1, j-1}^p + 1
    \]
    \textbf{Giustificazione:} L'aggiunta di un simbolo non rosso non modifica il numero di simboli rossi né il suo resto modulo $3$, incrementando di $1$ la soluzione di $(i-1, j-1, p)$.

    \item \textbf{Caso $x_i = y_j$ con $\text{col}(x_i) = \text{rosso}$:}
    Aggiungendo un simbolo rosso, il resto modulo $3$ avanza di $1$; pertanto per ottenere resto $p$ includendo $x_i = y_j$ occorre provenire dallo stato $(p - 1) \bmod 3$:
    \begin{itemize}
        \item Se $p = 0$:
        \[
            c_{i,j}^0 = \max \{ c_{i-1, j-1}^2 + 1, \, c_{i-1, j-1}^0 \}
        \]
        \textbf{Giustificazione:} Un multiplo di $3$ ($p=0$) si ottiene aggiungendo il rosso corrente a una sottosequenza con resto $2$ ($2 + 1 \equiv 0 \pmod 3$), oppure scartando il rosso corrente e mantenendo la soluzione di resto $0$.
        \item Se $p = 1$:
        \[
            c_{i,j}^1 = \max \{ c_{i-1, j-1}^0 + 1, \, c_{i-1, j-1}^1 \} = c_{i-1, j-1}^0 + 1
        \]
        \textbf{Giustificazione:} Un resto $1$ si ottiene aggiungendo il rosso corrente a una sottosequenza con resto $0$ ($0 + 1 \equiv 1 \pmod 3$).
        \item Se $p = 2$:
        \[
            c_{i,j}^2 = \max \{ c_{i-1, j-1}^1 + 1, \, c_{i-1, j-1}^2 \} = c_{i-1, j-1}^1 + 1
        \]
        \textbf{Giustificazione:} Un resto $2$ si ottiene aggiungendo il rosso corrente a una sottosequenza con resto $1$ ($1 + 1 \equiv 2 \pmod 3$).
    \end{itemize}
\end{itemize}

\subsection{Soluzione del problema principale}
La lunghezza della più lunga sottosequenza comune tra $X$ e $Y$ con un numero di simboli rossi multiplo di $3$ è:
\[
    \text{Soluzione} = c_{m, n}^0
\]

\section{Algoritmo Bottom-Up}
\begin{algorithm}[htbp]
\caption{Calcolo Bottom-Up per LCS con numero di rossi multiplo di 3}
\begin{algorithmic}[1]
\Procedure{LCSM3R}{$X, Y, \text{col}$}
    \State $m \gets |X|$
    \State $n \gets |Y|$
    \State $c \gets \text{nuova matrice } c[0 \dots m, \, 0 \dots n, \, 0 \dots 2]$
    \Comment{Caso base}
    \For{$i = 0 \textbf{ to } m$}
        \State $c[i, 0, 0] \gets 0$
        \For{$p = 1 \textbf{ to } 2$}
            \State $c[i, 0, p] \gets -\infty$
        \EndFor
    \EndFor
    \For{$j = 0 \textbf{ to } n$}
        \State $c[0, j, 0] \gets 0$
        \For{$p = 1 \textbf{ to } 2$}
            \State $c[0, j, p] \gets -\infty$
        \EndFor
    \EndFor
    \Comment{Passo ricorsivo}
    \For{$i = 1 \textbf{ to } m$}
        \For{$j = 1 \textbf{ to } n$}
            \For{$p = 0 \textbf{ to } 2$}
                \If{$x_i \neq y_j$}
                    \State $c[i, j, p] \gets \max \{ c[i-1, j, p], \, c[i, j-1, p] \}$
                \Else
                    \If{$\text{col}(x_i) \neq \text{rosso}$}
                        \State $c[i, j, p] \gets c[i-1, j-1, p] + 1$
                    \Else
                        \If{$p == 0$}
                            \State $c[i, j, 0] \gets \max \{ c[i-1, j-1, 2] + 1, \, c[i-1, j-1, 0] \}$
                        \ElsIf{$p == 1$}
                            \State $c[i, j, 1] \gets c[i-1, j-1, 0] + 1$
                        \Else
                            \State $c[i, j, 2] \gets c[i-1, j-1, 1] + 1$
                        \EndIf
                    \EndIf
                \EndIf
            \EndFor
        \EndFor
    \EndFor
    \State \Return $c$
\EndProcedure
\end{algorithmic}
\end{algorithm}

\section{Algoritmo di ricostruzione}
\begin{algorithm}[htbp]
\caption{Ricostruzione della sequenza LCS con numero di rossi multiplo di 3}
\begin{algorithmic}[1]
\Procedure{LCSM3RR}{$c, X, Y, i, j, p$}
    \If{$i == 0 \lor j == 0$}
        \State \Return
    \Else
        \If{$x_i \neq y_j$}
            \If{$c[i-1, j, p] \ge c[i, j-1, p]$}
                \State \Call{LCSM3RR}{$c, X, Y, i-1, j, p$}
            \Else
                \State \Call{LCSM3RR}{$c, X, Y, i, j-1, p$}
            \EndIf
        \Else
            \If{$\text{col}(x_i) \neq \text{rosso}$}
                \State \Call{LCSM3RR}{$c, X, Y, i-1, j-1, p$}
                \State \textbf{print} $x_i$
            \Else
                \If{$p == 0$}
                    \If{$c[i-1, j-1, 2] + 1 \ge c[i-1, j-1, 0]$}
                        \State \Call{LCSM3RR}{$c, X, Y, i-1, j-1, 2$}
                        \State \textbf{print} $x_i$
                    \Else
                        \State \Call{LCSM3RR}{$c, X, Y, i-1, j-1, 0$}
                    \EndIf
                \ElsIf{$p == 1$}
                    \State \Call{LCSM3RR}{$c, X, Y, i-1, j-1, 0$}
                    \State \textbf{print} $x_i$
                \Else
                    \State \Call{LCSM3RR}{$c, X, Y, i-1, j-1, 1$}
                    \State \textbf{print} $x_i$
                \EndIf
            \EndIf
        \EndIf
    \EndIf
\EndProcedure
\end{algorithmic}
\end{algorithm}

\noindent
\textbf{Chiamata principale (Driver):} Si invoca $\Call{LCSM3RR}{c, X, Y, m, n, 0}$.

\section{Analisi della complessità}
\begin{itemize}
    \item \textbf{Complessità spaziale:} $\Theta(3 \cdot m \cdot n) = \Theta(m \cdot n)$.
    \item \textbf{Complessità temporale (Bottom-Up):} $\Theta(3 \cdot m \cdot n) = \Theta(m \cdot n)$.
    \item \textbf{Complessità temporale (Ricostruzione):} $O(m + n)$.
\end{itemize}